\section{Perspectiva de ingeniería: flujo de trabajo de Q-Learning}

\subsection{Ciclo cerrado datos → entrenamiento → despliegue}

El ciclo cerrado típico de un proyecto de ingeniería incluye: recolección de datos → entrenamiento del modelo → revisión de la evaluación → Sim rollout + human feedback → monitoreo en producción.

\begin{importantbox}{Nodos clave del ciclo cerrado}
1. Recolectar transiciones de alta calidad, evitando que el agente haga trampa. 2. Usar PER + Double Q + regularización para estabilizar el entrenamiento. 3. Reforzar el paradigma de evaluación: replay con humano en el ciclo, etiquetado de casos de fallo. 4. Establecer límites de monitoreo: error de TD, varianza de la recompensa, latencia.
\end{importantbox}

\subsection{Estrategia de monitoreo en producción}

Después del despliegue, el monitoreo se centra en:

\begin{itemize}
    \item deriva del `TD error` (si aumenta, es necesario reentrenar o hacer rollback)
    \item varianza de la recompensa (si disminuye, puede indicar un colapso de la conducta)
    \item `episode success rate` y `reward per step`, para asegurar que el desempeño no se degrade
\end{itemize}

\subsection{Resumen de la sección}

El despliegue de ingeniería enfatiza la higiene del replay buffer, la vigilancia del error de TD y el rollback impulsado por telemetría.

